% BEGIN WOTSC_SECURITY_PROOF_INTEGRATION_V1A

\subsection{Source-bound security model for the experimental WOTS+C profile}
\label{subsec:wotsc_security_model}

Beyond the implementation-level correctness evidence above, we developed a
source-bound security reduction for the canonical experimental
XMSS$^{\mathrm{MT}}$ profile using WOTS+C on both height-10 layers. The
analysis is deliberately separated into a primary standard-assumption theorem
and an optional numerical corollary under explicitly idealized primitive and
state-management assumptions. In particular, the implementation measurements
are not used as substitutes for adversarial query bounds.

Let $q_{F,0}$ and $q_{F,1}$ denote the adversary's addressed cut-oracle query
budgets in the two layer-local WOTS+C reductions. These quantities are
security-game resource parameters; they are not the number of honest
implementation calls to the chain function and are not the complete
implementation-lifetime hash-call counts. For either layer $j\in\{0,1\}$,
the local reduction gives
\[
  \operatorname{Adv}_{\mathrm{WOTSC},j}(q_{F,j})
  \le
  \frac{8(4q_{F,j}+961)^2}{2^{256}}.
\]
Here, the additive constant $961=960+1$ arises from the fixed verification
overhead in the reconstructed marking-oracle reduction together with the
residual-search lemma's final $+1$ allowance.

For the composed two-layer construction, the primary forgery bound is
\[
\begin{aligned}
\Pr[\mathrm{Forge}_{\mathrm{XMSSMT}}]
\le{}&
  \operatorname{Adv}^{\mathrm{prf}}_{\mathrm{PRF\mbox{-}keygen}}
  + 2^{-256}
  + \operatorname{Adv}_{H_{\mathrm{msg}}\text{-}\mathrm{bind}}
  + \Pr[E_{\mathrm{state}}] \\
 &+
  \operatorname{Adv}_{\mathrm{WOTSC},0}(q_{F,0})
  + \operatorname{Adv}_{\mathrm{WOTSC},1}(q_{F,1})
  + \operatorname{Adv}_{H_1\text{-}\mathrm{bind}}
  + \Pr[E_{\mathrm{grind}}].
\end{aligned}
\]
The $2^{-256}$ term corresponds to equality of independently and uniformly
generated \texttt{SK\_PRF} and \texttt{PUB\_SEED}. The shared
\texttt{SK\_SEED} reduction is charged once globally rather than once per
XMSS$^{\mathrm{MT}}$ layer. The source-bound resource analysis limits the
PRF-keygen reduction to at most $67{,}174{,}400$ distinct challenge inputs,
the global DOMAIN-1 ledger to at most $334{,}972{,}196$ queries, and the
complete $H_{\mathrm{msg}}$ ledger used by the theorem to at most
$1{,}048{,}577$ queries. The corresponding PRF and binding advantages remain
symbolic in the primary theorem rather than being assigned numerical values
solely from the 256-bit SHA-256 output length.

\subsection{State and grinding assumptions}
\label{subsec:wotsc_state_grinding}

For uninterrupted sequential execution, the audited implementation advances
the XMSS$^{\mathrm{MT}}$ signing index without legitimate index reuse.
Deployment-level reuse remains a separate event, denoted
$E_{\mathrm{state}}$, because process rollback, restoration of stale private
state, cloned signer images, concurrent writers, or non-durable index
retirement can invalidate the stateful-signature assumption. Accordingly,
$\Pr[E_{\mathrm{state}}]$ remains symbolic in the primary theorem.

WOTS+C searches a 32-bit counter space for a digest whose 64 hexadecimal
digits have the required fixed sum of 480. The successful counter is
deterministically bounded by
$0\le c\le 2^{32}-1$. Under the separate model in which distinct canonical
type-7 counter inputs produce independent uniform 256-bit outputs, the
single-trial acceptance probability is approximately
$0.0107922589805859367$, and exhaustion of the complete $2^{32}$ candidate
space is bounded below the probability $2^{-67{,}235{,}846}$. This
grinding-exhaustion statement is conditional on that explicit ideal-uniform
model; it is not asserted as an unconditional concrete property of SHA-256.

\subsection{Conditional numerical corollary}
\label{subsec:wotsc_conditional_corollary}

For numerical interpretation only, consider an additional idealized model in
which the PRF-keygen challenge interface is an ideal random function,
distinct canonical $H_{\mathrm{msg}}$ inputs yield independent uniform
256-bit outputs, and distinct canonical DOMAIN-1 inputs likewise yield
independent uniform 256-bit outputs. Further assume a single authorized
signer with atomic durable index retirement, no rollback, no stale-state
restoration, no signer-state cloning, and no concurrent signing from the same
private-state snapshot.

Under these assumptions, the binding terms can be bounded by the birthday
expressions
\[
  \frac{\binom{1{,}048{,}577}{2}}{2^{256}}
  \quad\text{and}\quad
  \frac{\binom{334{,}972{,}196}{2}}{2^{256}},
\]
corresponding to approximately $217.00$ and $200.36$ bits, respectively.
If, in addition,
\[
 q_{F,0},q_{F,1}
 \le
 576{,}460{,}752{,}303{,}423{,}247,
\]
then each layer-local WOTS+C term is at most $2^{-131}$. Combining these
terms with the key-equality and grinding-exhaustion terms gives the frozen
conditional aggregate bound
\[
  \Pr[\mathrm{Forge}_{\mathrm{XMSSMT}}]
  \le
  7.3468396926392969\times 10^{-40}
  < 2^{-128}.
\]
The corresponding aggregate exponent is approximately $130$ bits.
Consequently, the idealized corollary supports at least a 128-bit security
target under the stated primitive, adversary-resource, and durable-state
assumptions.

This numerical corollary must not be interpreted as an unconditional
128-bit or 130-bit concrete security proof for the real SHA-256
implementation. In particular, the analysis does not assign a zero concrete
PRF advantage to the implemented SHA-256 construction, nor does it claim
that the real $H_{\mathrm{msg}}$ and DOMAIN-1 binding advantages are exactly
the idealized birthday values above. The primary manuscript claim therefore
remains the parameterized standard-assumption theorem.

% END WOTSC_SECURITY_PROOF_INTEGRATION_V1A

